\documentclass[12pt,a4paper]{article}

\usepackage[utf8]{inputenc}
\usepackage[T2A]{fontenc}
\usepackage[russian]{babel}
\usepackage{amsmath,amssymb}
\usepackage{graphicx}
\usepackage{booktabs}
\usepackage{pgfplots}
\pgfplotsset{compat=1.18}
\usepackage{geometry}
\geometry{top=2cm,bottom=2cm,left=2.5cm,right=1.5cm}
\usepackage{float}
\usepackage{makecell}

\title{\textbf{Лабораторная работа 1.1.1}\\
Определение удельного сопротивления проводника}

\author{
Студент: \underline{Шеховцов Василий Максимович}\\
Группа: \underline{Б03-601}\\
Преподаватель: \underline{Ивлев Д. А.}
}

\date{Московский физико-технический институт\\
(национальный исследовательский университет)\\
9 сентября 2026 г.\\}

\begin{document}

\maketitle

\section*{Аннотация}

В работе измеряется удельное сопротивление тонкой нихромовой проволоки
круглого сечения. Сопротивление проволоки определяется двумя независимыми
методами: 1) по угловому коэффициенту зависимости напряжения на проволоке
от протекающего через неё тока, измеряемых аналоговым вольтметром и
цифровым амперметром; 2) с помощью моста постоянного тока Р4833.
Геометрические размеры образца (диаметр и длина) измеряются линейкой,
штангенциркулем и микрометром. Для каждого метода проводится оценка
случайных и систематических погрешностей, после чего результаты двух
методов сравниваются между собой и с табличными значениями удельного
сопротивления нихрома.

\section*{Цель работы}

Определить удельное электрическое сопротивление нихромовой проволоки,
используя амперметр, вольтметр и мост постоянного тока. Оценить погрешности измерений.

\section*{Оборудование}

\begin{itemize}
    \item Нихромовая проволока на рейке с подвижным контактом;
    \item микрометр (цена деления $0.01$ мм);
    \item штангенциркуль;
    \item линейка;
    \item аналоговый вольтметр;
    \item цифровой амперметр GDM-8245;
    \item реостат;
    \item источник ЭДС;
    \item мост постоянного тока Р4833;
    \item соединительные провода, ключ.
\end{itemize}

\section*{1. Теоретические сведения}

Сопротивление однородного проводника круглого сечения связано с удельным
сопротивлением $\rho$ формулой

\begin{equation}
R=\rho\frac{l}{S},
\qquad
\rho=\frac{RS}{l}
=\frac{R\pi d^2}{4l},
\label{eq:rho}
\end{equation}

где $l$ --- длина проводника, $d$ --- его диаметр,
$S=\pi d^2/4$ --- площадь поперечного сечения.

Согласно закону Ома:

\begin{equation}
U=RI.
\label{eq:ohm}
\end{equation}

Для измерения напряжения и тока использовалась схема, изображённая на
рис.1 \\ амперметр измеряет полный ток $I_A$, а вольтметр
подключён параллельно исследуемой проволоке.

\begin{figure}[H]
\centering
\includegraphics[width=0.55\textwidth]{C:/Users/basil/OneDrive/Изображения/Screenshots/Снимок экрана 2026-09-12 163200.png}
\caption{Схема измерения вольт-амперной характеристики проволоки}
\label{fig:scheme}
\end{figure}

Ввиду неидеальности вольтметра необходимо учесть его конечное внутреннее 
сопротивление $R_V$. Показания амперметра $I_A$ и вольтметра $V_B$ связаны соотношением 

\begin{equation}
V_B=R'I_A,
\end{equation}

где $R'$ --- сопротивление параллельно соединённых проволоки и вольтметра:

\begin{equation}
\frac{1}{R'}=\frac{1}{R}+\frac{1}{R_V}.
\end{equation}

При $R_V\gg R,R'$ график зависимости $V_B(I_A)$ должен представлять
прямую, угловой коэффициент которой равен $R'$.

Истинное сопротивление проволоки находится с учётом поправки на конечное
сопротивление вольтметра:

\begin{equation}
R=\frac{R_VR'}{R_V-R'}
\approx
R'\left(1+\frac{R'}{R_V}\right).
\label{eq:correction}
\end{equation}

Для определения сопротивления по серии измерений используется метод
наименьших квадратов. Если аппроксимирующая прямая проходит через начало
координат,

\begin{equation}
U=RI,
\end{equation}

то

\begin{equation}
R=
\frac{\langle UI\rangle}
{\langle I^2\rangle}.
\label{eq:mnk}
\end{equation}

Случайная погрешность углового коэффициента определяется как

\begin{equation}
\sigma_{\text{сл}}=
\sqrt{
\frac{1}{n-1}
\left(
\frac{\langle U^2\rangle}
{\langle I^2\rangle}
-R^2
\right)
}.
\label{eq:random}
\end{equation}

Полная погрешность сопротивления:

\begin{equation}
\sigma_R=
\sqrt{
\left(\sigma_{\text{сл}}\right)^2+
\left(\sigma_{\text{сист}}\right)^2
}.
\label{eq:fullR}
\end{equation}


\section*{2. Результаты прямых измерений}

\subsection*{2.1. Измерение диаметра проволоки}

Диаметр проволоки измерялся в 10 различных точках сначала
штангенциркулем, затем микрометром.

\begin{table}[H]
\centering
\caption{Результаты измерения диаметра проволоки}
\label{tab:diam}
\begin{tabular}{c|c|c}
\toprule
№ измерения & Штангенциркуль, мм & Микрометр, мм\\
\midrule
1 & 0.40 & 0.36\\
2 & 0.40 & 0.37\\
3 & 0.40 & 0.35\\
4 & 0.40 & 0.35\\
5 & 0.40 & 0.34\\
6 & 0.40 & 0.36\\
7 & 0.40 & 0.35\\
8 & 0.40 & 0.34\\
9 & 0.40 & 0.35\\
10 & 0.40 & 0.35\\
\bottomrule
\end{tabular}
\end{table}

Среднее значение диаметра по микрометру:

\begin{equation}
\bar d=
\frac{1}{10}\sum_{i=1}^{10}d_i
=
0.352\ \text{мм}.
\end{equation}

Среднее значение по штангенциркулю:

\begin{equation}
\bar d_{\text{шт}}
=
0.4\ \text{мм}.
\end{equation}

Стандартное отклонение результатов микрометра:

\begin{equation}
s_d=
\sqrt{
\frac{1}{n-1}
\sum_{i=1}^{n}(d_i-\bar d)^2
}
=
0.00919\ \text{мм}.
\end{equation}

Случайная погрешность среднего:

\begin{equation}
\sigma_{\bar d}
=
\frac{s_d}{\sqrt n}
=
0.00291\ \text{мм}.
\end{equation}

Для микрометра принимаем системную погрешность

\begin{equation}
\Delta d_{\text{мкм}}
=
0.01\ \text{мм}.
\end{equation}

Полная погрешность диаметра:

\begin{equation}
\sigma_d=
\sqrt{
\sigma_{\bar d}^2+
\Delta d_{\text{мкм}}^2
}
=
0.0104\ \text{мм}.
\end{equation}

Таким образом,

\begin{equation}
\boxed{
d=(0.352\pm0.010)\ \text{мм}
}.
\end{equation}

Относительная погрешность диаметра составляет

\begin{equation}
\varepsilon_d=
\frac{\sigma_d}{d}
\approx
2.96\%.
\end{equation}

Заметим, что штангенциркуль даёт систематически завышенное значение
диаметра ($0.4$ мм против $0.352$ мм по микрометру) и при этом нулевой
разброс по всем 10 измерениям. Это объясняется тем, что цена деления
штангенциркуля ($0.05$ мм) значительно грубее диаметра тонкой проволоки,
поэтому все отсчёты попадают в одно и то же деление шкалы.


\subsection*{2.2. Расчёт площади поперечного сечения}

Площадь поперечного сечения проволоки:

\begin{equation}
S=
\frac{\pi d^2}{4}
=
\frac{\pi(0.352\cdot10^{-3})^2}{4}
=
9.73\cdot10^{-8}\ \text{м}^2.
\end{equation}

Относительная погрешность площади:

\begin{equation}
\frac{\sigma_S}{S}
=
2\frac{\sigma_d}{d}
=
2\cdot\frac{0.0104}{0.352}
\approx
0.0592
=
5.92\%.
\end{equation}

Абсолютная погрешность:

\begin{equation}
\sigma_S=
0.0592\cdot9.73\cdot10^{-8}
\approx
0.58\cdot10^{-8}\ \text{м}^2.
\end{equation}

Итого:

\begin{equation}
\boxed{
S=(9.73\pm0.58)\cdot10^{-8}\ \text{м}^2
}.
\end{equation}


\subsection*{2.3. Характеристики измерительных приборов}

Для вольтметра использованы характеристики из паспорта/теории
предыдущей лабораторной работы: класс точности $0.5$, шкала 150 делений,
предел $0.75$ В, цена деления $5$ мВ, чувствительность $200$ дел./В
и внутреннее сопротивление $5$ кОм.


\begin{table}[H]
\centering
\caption{Основные характеристики вольтметра и амперметра}
\label{tab:instruments}
\begin{tabular}{c|c|c}
\toprule
Параметр & Вольтметр & Амперметр GDM-8245\\
\midrule
Тип & аналоговый & цифровой\\
Класс точности & 0.5 & 0.5\\
Предел измерений & 0.75 В & 500 мА\\
Число делений & 150 & ---\\
Цена деления & 0.005 В & ---\\
Чувствительность & 200 дел./В & ---\\
Абсолютная погрешность &
$\pm0.00375$ В &
$\pm(0.002X+2k)$\\
Внутреннее сопротивление & 5 кОм & 1.4 Ом\\
\bottomrule
\end{tabular}
\end{table}

Для амперметра цифрового типа цена деления и чувствительность
в дел./ед. величины не определяются. Согласно паспорту:

\begin{equation}
\Delta I=
\pm(0.002I+2k),
\end{equation}

где $I$ выражено в мА, а $k=0.01$ мА. Следовательно,

\begin{equation}
\Delta I=
\pm(0.002I+0.02)\ \text{мА}.
\end{equation}

Для вольтметра максимальная системная погрешность по классу
точности:

\begin{equation}
\Delta U=
0.005\cdot0.75
=
0.00375\ \text{В}.
\end{equation}


\section*{3. Измерение сопротивления с помощью амперметра и вольтметра}

Цена деления вольтметра составляет $0.005$ В, поэтому напряжение
для каждой точки рассчитывается как

\begin{equation}
U=N_V\cdot0.005\ \text{В},
\end{equation}

где $N_V$ --- показание шкалы в делениях.

Ниже приведены исходные показания и рассчитанные напряжения.


\subsection*{3.1. Длина $l=50$ см}

\begin{table}[H]
\centering
\caption{Измерения при длине проволоки 50 см}
\begin{tabular}{c|c|c|c}
\toprule
№ & $V$, дел & $V$, В & $I$, мА\\
\midrule
1&34.5&0.1725&33.4\\
2&36.1&0.1805&35.1\\
3&40.0&0.2000&38.5\\
4&44.0&0.2200&42.4\\
5&51.5&0.2575&50.1\\
6&86.7&0.4335&83.2\\
7&79.0&0.3950&78.9\\
8&71.5&0.3575&68.9\\
9&65.5&0.3275&63.4\\
10&58.0&0.2900&56.2\\
11&55.0&0.2750&53.4\\
12&50.0&0.2500&48.6\\
\bottomrule
\end{tabular}
\end{table}

По МНК угловой коэффициент графика даёт измеренное сопротивление
параллельной комбинации $R'$:

\begin{equation}
R'_{50}=5.148\ \text{Ом}.
\end{equation}

С учётом конечного сопротивления вольтметра $R_V=5000$ Ом:

\begin{equation}
R_{50}
=
\frac{5000\cdot5.148}{5000-5.148}
=
5.153\ \text{Ом}.
\end{equation}

Случайная погрешность:

\begin{equation}
\sigma_{\text{сл}}
=
0.020\ \text{Ом}.
\end{equation}

Для максимального тока $I_{\max}=83.2$ мА:

\begin{equation}
\Delta I_{\max}
=
0.002\cdot83.2+0.02
=
0.1864\ \text{мА}.
\end{equation}

Систематическая погрешность:

\begin{equation}
\sigma_{\text{сист}}
=
R
\sqrt{
\left(
\frac{\Delta U}{U_{\max}}
\right)^2+
\left(
\frac{\Delta I}{I_{\max}}
\right)^2
}
\approx
0.046\ \text{Ом}.
\end{equation}

Полная погрешность:

\begin{equation}
\sigma_R=
\sqrt{
0.020^2+0.046^2
}
\approx
0.050\ \text{Ом}.
\end{equation}

Итого:

\begin{equation}
\boxed{
R_{50}=(5.15\pm0.05)\ \text{Ом}
}.
\end{equation}


\subsection*{3.2. Длина $l=30$ см}

\begin{table}[H]
\centering
\caption{Измерения при длине проволоки 30 см}
\begin{tabular}{c|c|c|c}
\toprule
№ & $V$, дел & $V$, В & $I$, мА\\
\midrule
1&17.0&0.0850&27.4\\
2&26.0&0.1300&42.7\\
3&25.1&0.1255&47.1\\
4&40.0&0.2000&64.5\\
5&46.3&0.2315&75.4\\
6&52.5&0.2625&86.5\\
7&48.5&0.2425&79.2\\
8&44.0&0.2200&71.8\\
9&41.2&0.2060&67.5\\
10&39.0&0.1950&63.1\\
11&36.0&0.1800&59.1\\
12&32.0&0.1600&52.1\\
\bottomrule
\end{tabular}
\end{table}

По МНК получаем

\begin{equation}
R'_{30}=3.045\ \text{Ом}.
\end{equation}

С учётом поправки:

\begin{equation}
R_{30}
=
\frac{5000\cdot3.045}{5000-3.045}
=
3.047\ \text{Ом}.
\end{equation}

Получаем

\begin{equation}
\sigma_{\text{сл}}
=
0.026\ \text{Ом},
\qquad
\sigma_{\text{сист}}
=
0.044\ \text{Ом},
\end{equation}

и

\begin{equation}
\boxed{
R_{30}=(3.04\pm0.05)\ \text{Ом}
}.
\end{equation}


\subsection*{3.3. Длина $l=20$ см}

\begin{table}[H]
\centering
\caption{Измерения при длине проволоки 20 см}
\begin{tabular}{c|c|c|c}
\toprule
№ & $V$, дел & $V$, В & $I$, мА\\
\midrule
1&31.7&0.1585&75.3\\
2&34.0&0.1700&81.1\\
3&36.0&0.1800&86.2\\
4&41.2&0.2060&98.6\\
5&58.0&0.2900&114.8\\
6&68.3&0.3415&165.0\\
7&58.0&0.2900&137.9\\
8&51.0&0.2550&121.7\\
9&46.0&0.2300&111.6\\
10&40.5&0.2025&97.6\\
11&36.0&0.1800&86.9\\
12&31.5&0.1575&75.2\\
\bottomrule
\end{tabular}
\end{table}

По МНК получаем

\begin{equation}
R'_{20}=2.126\ \text{Ом}.
\end{equation}

С учётом поправки:

\begin{equation}
R_{20}
=
\frac{5000\cdot2.126}{5000-2.126}
=
2.127\ \text{Ом}.
\end{equation}

Получаем

\begin{equation}
\sigma_{\text{сл}}
=
0.039\ \text{Ом},
\qquad
\sigma_{\text{сист}}
=
0.024\ \text{Ом},
\end{equation}

и

\begin{equation}
\boxed{
R_{20}=(2.13\pm0.05)\ \text{Ом}
}.
\end{equation}


\section*{4. Общий график зависимости $V=f(I)$}

На одном графике представлены экспериментальные зависимости
$V(I)$ для всех трёх длин проволоки и соответствующие прямые,
полученные методом наименьших квадратов.

\begin{figure}[H]
\centering

\begin{tikzpicture}
\begin{axis}[
    width=0.95\textwidth,
    height=0.62\textwidth,
    xlabel={$I$, мА},
    ylabel={$V$, В},
    xmin=0,
    xmax=180,
    ymin=0,
    ymax=0.48,
    grid=both,
    legend style={
        at={(0.98,0.02)},
        anchor=south east,
        font=\small,
        fill=white,
        draw=black
    },
]

% =====================================
% l = 20 см — эксперимент
% =====================================

\addplot[
    only marks,
    mark=*,
    mark size=1.6pt
]
coordinates {
    (75.3,0.1585)
    (81.1,0.1700)
    (86.2,0.1800)
    (98.6,0.2060)
    (114.8,0.2900)
    (165.0,0.3415)
    (137.9,0.2900)
    (121.7,0.2550)
    (111.6,0.2300)
    (97.6,0.2025)
    (86.9,0.1800)
    (75.2,0.1575)
};
\addlegendentry{$l=20$ см, эксперимент}

% l = 20 см — МНК
\addplot[
    domain=0:180,
    samples=2,
    thick
]
{0.00212603*x};
\addlegendentry{$l=20$ см, МНК: $R'=2.126\,\Omega$}


% =====================================
% l = 30 см — эксперимент
% =====================================

\addplot[
    only marks,
    mark=square*,
    mark size=1.6pt
]
coordinates {
    (27.4,0.0850)
    (42.7,0.1300)
    (47.1,0.1255)
    (64.5,0.2000)
    (75.4,0.2315)
    (86.5,0.2625)
    (79.2,0.2425)
    (71.8,0.2200)
    (67.5,0.2060)
    (63.1,0.1950)
    (59.1,0.1800)
    (52.1,0.1600)
};
\addlegendentry{$l=30$ см, эксперимент}

% l = 30 см — МНК
\addplot[
    domain=0:180,
    samples=2,
    thick
]
{0.00304470*x};
\addlegendentry{$l=30$ см, МНК: $R'=3.045\,\Omega$}


% =====================================
% l = 50 см — эксперимент
% =====================================

\addplot[
    only marks,
    mark=triangle*,
    mark size=1.8pt
]
coordinates {
    (33.4,0.1725)
    (35.1,0.1805)
    (38.5,0.2000)
    (42.4,0.2200)
    (50.1,0.2575)
    (83.2,0.4335)
    (78.9,0.3950)
    (68.9,0.3575)
    (63.4,0.3275)
    (56.2,0.2900)
    (53.4,0.2750)
    (48.6,0.2500)
};
\addlegendentry{$l=50$ см, эксперимент}

% l = 50 см — МНК
\addplot[
    domain=0:180,
    samples=2,
    thick
]
{0.00514751*x};
\addlegendentry{$l=50$ см, МНК: $R'=5.148\,\Omega$}

\end{axis}
\end{tikzpicture}

\caption{Зависимость напряжения от силы тока для трёх длин
нихромовой проволоки}
\label{fig:VI_all}

\end{figure}


\section*{5. Измерение сопротивления мостом Р4833}

По данным прямых измерений с помощью моста получены:

\begin{table}[H]
\centering
\caption{Результаты измерений сопротивления мостом}
\label{tab:bridge}
\begin{tabular}{c|c|c}
\toprule
Длина $l$, см & $R_{\text{мост}}$, Ом & Погрешность, Ом\\
\midrule
20 & 2.1299 & $\pm0.010$\\
30 & 3.0130 & $\pm0.010$\\
50 & 5.1789 & $\pm0.010$\\
\bottomrule
\end{tabular}
\end{table}


\subsection*{5.1. Сводная таблица результатов}

В таблице приведены результаты, полученные при различных длинах
проволоки. $R_0$ --- сопротивление, измеренное мостом Р4833,
$R_{\text{ср}}$ --- сопротивление, полученное методом наименьших
квадратов до внесения поправки, $R_{\text{пр}}$ --- сопротивление
с учётом поправки на конечное сопротивление вольтметра,
$\sigma_{\text{сл}}$ и $\sigma_{\text{сист}}$ --- случайная и
систематическая составляющие погрешности, $\sigma_R$ --- полная
погрешность.

\begin{table}[H]
\centering
\caption{Результаты измерения сопротивления проволоки}
\label{tab:summary}
\begin{tabular}{|c|c|c|}
\hline
$l=20$ см & $l=30$ см & $l=50$ см \\
\hline
\makecell{$R_0=2.1299$ Ом \\ (по Р4833)} &
\makecell{$R_0=3.0130$ Ом \\ (по Р4833)} &
\makecell{$R_0=5.1789$ Ом \\ (по Р4833)} \\
\hline
$R_{\text{ср}}=2.126$ Ом & $R_{\text{ср}}=3.045$ Ом & $R_{\text{ср}}=5.148$ Ом \\
\hline
$R_{\text{пр}}=2.127$ Ом & $R_{\text{пр}}=3.047$ Ом & $R_{\text{пр}}=5.153$ Ом \\
\hline
$\sigma_{\text{сл}}=0.039$ Ом & $\sigma_{\text{сл}}=0.026$ Ом & $\sigma_{\text{сл}}=0.020$ Ом \\
\hline
$\sigma_{\text{сист}}=0.024$ Ом & $\sigma_{\text{сист}}=0.044$ Ом & $\sigma_{\text{сист}}=0.046$ Ом \\
\hline
$\sigma_R=0.046$ Ом & $\sigma_R=0.051$ Ом & $\sigma_R=0.050$ Ом \\
\hline
\end{tabular}
\end{table}


Сравним результаты двух методов:

\begin{table}[H]
\centering
\caption{Сравнение сопротивлений, полученных двумя методами}
\begin{tabular}{c|c|c|c}
\toprule
$l$, см & $R_{\text{АВ}}$, Ом & $R_{\text{мост}}$, Ом & Разность, Ом\\
\midrule
20 & 2.127 & 2.1299 & $-0.0029$\\
30 & 3.047 & 3.0130 & $+0.0340$\\
50 & 5.153 & 5.1789 & $-0.0259$\\
\bottomrule
\end{tabular}
\end{table}


\section*{6. Определение удельного сопротивления}

Для каждой длины используем результат, полученный с помощью моста:

\begin{equation}
\rho=
\frac{R\pi d^2}{4l}.
\end{equation}

При $l=20$ см:

\begin{equation}
\rho_{20}
=
1.036\cdot10^{-6}\ \text{Ом}\cdot\text{м}.
\end{equation}

При $l=30$ см:

\begin{equation}
\rho_{30}
=
0.977\cdot10^{-6}\ \text{Ом}\cdot\text{м}.
\end{equation}

При $l=50$ см:

\begin{equation}
\rho_{50}
=
1.008\cdot10^{-6}\ \text{Ом}\cdot\text{м}.
\end{equation}

Среднее значение:

\begin{equation}
\bar\rho=
\frac{
\rho_{20}+\rho_{30}+\rho_{50}
}{3}
=
1.007\cdot10^{-6}\ \text{Ом}\cdot\text{м}.
\end{equation}

Основной вклад в погрешность даёт измерение диаметра. Для каждого опыта:

\begin{equation}
\left(
\frac{\sigma_\rho}{\rho}
\right)^2
=
\left(
2\frac{\sigma_d}{d}
\right)^2
+
\left(
\frac{\sigma_R}{R}
\right)^2
+
\left(
\frac{\sigma_l}{l}
\right)^2.
\end{equation}

Для трёх измерений получаем примерно:

\begin{align}
\rho_{20}
&=
(1.036\pm0.062)
\cdot10^{-6}\ \text{Ом}\cdot\text{м},
\\
\rho_{30}
&=
(0.977\pm0.058)
\cdot10^{-6}\ \text{Ом}\cdot\text{м},
\\
\rho_{50}
&=
(1.008\pm0.060)
\cdot10^{-6}\ \text{Ом}\cdot\text{м}.
\end{align}

Поэтому окончательно:

\begin{equation}
\boxed{
\rho=
(1.01\pm0.06)
\cdot10^{-6}\ \text{Ом}\cdot\text{м}
}.
\end{equation}

Относительная погрешность результата составляет около $6\%$.

\section*{7. Обсуждение результатов и выводы}

В работе получено значение удельного сопротивления образца проволоки
из нихромового сплава с точностью около $6\%$. Табличные значения
для нихрома лежат в диапазоне
\[
\rho_{\text{табл}}=0{,}97\ldots1{,}14\cdot10^{-6}\ \text{Ом}\cdot\text{м}
\]
в зависимости от состава. Измеренное значение
\[
\rho=(1{,}01\pm0{,}06)\cdot10^{-6}\ \text{Ом}\cdot\text{м}
\]
попадает в этот диапазон, что подтверждает корректность проведённых
измерений.

Использованный метод измерения сопротивления позволил получить значения
$R$ с точностью порядка $1$---$2\%$, которая в основном ограничивалась
погрешностью аналогового вольтметра. Сравнение с результатами измерений
с помощью моста Р4833 показало совпадение результатов в пределах
погрешностей.

Точность измерения удельного сопротивления ограничивалась
измерением диаметра проволоки. Поскольку $\rho\propto d^2$, погрешность
диаметра оказывает значительное влияние на итоговый результат.
Увеличение числа измерений микрометром не позволяет повысить
точность, так как случайный разброс оказался меньше цены деления прибора.

\end{document}
